% TikZ 图：核裂变与核聚变模型
% Day 27：近代物理
\begin{tikzpicture}[scale=0.95,>=Latex]
  \tikzset{
    nucleus/.style={circle,draw=black,thick,fill=gray!12,minimum size=0.62cm},
    neutron/.style={circle,draw=blue!70!black,fill=blue!18,minimum size=0.24cm,inner sep=0pt},
    ray/.style={->,thick},
    energy/.style={->,very thick,orange!90!black},
    note/.style={font=\small}
  }

  % 裂变
  \begin{scope}
    \node[font=\bfseries] at (2.7,1.85) {核裂变：重核分裂};
    \node[neutron,label=below:{\scriptsize 中子}] (n0) at (0.2,0) {};
    \node[nucleus,minimum size=1.05cm] (u) at (1.55,0) {$^{235}\mathrm{U}$};
    \draw[ray,blue!70!black] (n0) -- (u);

    \node[nucleus] (a) at (3.25,0.55) {碎片};
    \node[nucleus] (b) at (3.25,-0.55) {碎片};
    \draw[ray] (u) -- (a);
    \draw[ray] (u) -- (b);

    \foreach \x/\y in {4.35/0.80,4.55/0.10,4.20/-0.65} {
      \node[neutron] at (\x,\y) {};
      \draw[ray,blue!70!black] (3.62,0) -- (\x,\y);
    }
    \draw[energy] (1.55,0.58) -- (0.90,1.25);
    \node[orange!90!black,font=\small] at (0.72,1.52) {能量};

    \node[note,align=center] at (2.65,-1.55) {中子诱发重核分裂\\新中子可能继续引发链式反应};
  \end{scope}

  % 分隔线
  \draw[gray!55,dashed] (5.15,-1.85) -- (5.15,2.05);

  % 聚变
  \begin{scope}[xshift=5.8cm]
    \node[font=\bfseries] at (2.35,1.85) {核聚变：轻核结合};
    \node[nucleus] (d) at (0.55,0.45) {$^2\mathrm{H}$};
    \node[nucleus] (t) at (0.55,-0.45) {$^3\mathrm{H}$};
    \draw[ray,red] (d) -- (1.45,0.12);
    \draw[ray,red] (t) -- (1.45,-0.12);

    \node[nucleus,minimum size=0.9cm] (he) at (2.35,0) {$^4\mathrm{He}$};
    \node[neutron,label=right:{\scriptsize 中子}] (n1) at (3.55,0.55) {};
    \draw[ray] (he) -- (n1);
    \draw[energy] (2.95,-0.55) -- (3.75,-0.95);
    \node[orange!90!black,right] at (3.55,-0.86) {能量};

    \node[note,align=center] at (2.35,-1.55) {太阳能量主要来自聚变\\需足够温度、密度和约束时间};
  \end{scope}
\end{tikzpicture}
